\title{DeepSeekMath: Pushing the Limits of Mathematical Reasoning in Open Language Models}

\begin{abstract}
Language models often struggle with mathematical reasoning due to its inherent structure and complexity. This work presents DeepSeekMath~7B, which extends the pre-training of DeepSeek-Coder-Base-v1.5 7B by incorporating 120B math-related tokens curated from Common Crawl, along with text and programming data. DeepSeekMath~7B attains a notable 51.7\% accuracy on the competition-level MATH benchmark, all without the aid of external toolkits or ensemble methods, nearing the capabilities of Gemini-Ultra and GPT-4. Evaluating self-consistency across 64 samples results in a 60.9\% score on MATH. DeepSeekMath~demonstrates its mathematical reasoning strength owing to two main advancements: Firstly, a carefully constructed pipeline for dataset selection leverages the breadth of open web resources. Secondly, we propose Group Relative Policy Optimization (GRPO), a novel adaptation of Proximal Policy Optimization (PPO), which not only improves the model's ability in mathematical problem solving but also makes PPO more memory-efficient.
\end{abstract}

\section{Introduction}

Large language models (LLMs) have transformed mathematical reasoning within artificial intelligence, driving notable progress on benchmarks such as the quantitative reasoning benchmark \citep{MATH} and the geometry reasoning benchmark \citep{trinh2024solving}. In addition, these systems have demonstrated their value in aiding humans with the resolution of intricate mathematical tasks \citep{tao}. Nevertheless, state-of-the-art models like GPT-4 \citep{gpt4} and Gemini-Ultra \citep{gemini} remain closed to public access, while open-source alternatives that are currently available exhibit substantially inferior performance.

This work presents DeepSeekMath, a specialized language model for mathematics that markedly surpasses other open-source alternatives and closely rivals the performance of GPT-4 on standard academic tests. Central to this advancement is the construction of the DeepSeekMath Corpus: a massive, high-caliber pre-training dataset containing 120B math-related tokens. To build this corpus, we leverage the Common Crawl (CC), employing a fastText-based classifier \citep{joulin2016fasttext}. In the classifier’s first iteration, positive examples are sourced from OpenWebMath \citep{openwebmath}, while an array of other web content forms the set of negative examples. The trained classifier is then used to identify further positive samples within the CC data, with subsequent rounds of manual annotation refining these selections. The classifier is subsequently retrained on the improved dataset, enhancing its extraction abilities. Our experimental assessments demonstrate the corpus’ outstanding quality: the DeepSeekMath-Base 7B model achieves 64.2\% on GSM8K \citep{gsm8k} and 36.2\% on the challenging MATH benchmark \citep{MATH}, exceeding the results obtained by Minerva 540B \citep{minerva}. Furthermore, thanks to the multilingual character of the DeepSeekMath Corpus, we observe notable progress in Chinese mathematics evaluations \citep{wei2023cmath,agieval}. We anticipate that our methodological contributions to mathematical data curation will serve as a valuable foundation for future research and recognize there remains substantial potential for further refinement.

DeepSeekMath-Base builds upon DeepSeek-Coder-Base-v1.5 7B \citep{deepseek-coder}, given our observation that initializing from a model trained on code yields better results than starting with a general-purpose LLM. Additionally, our findings reveal that mathematical training leads to improved performance on MMLU \citep{mmlu} and BBH benchmarks \citep{bbh}, suggesting enhancements not just in mathematical reasoning but also in the model's overall general reasoning skills.

Following pre-training, we conduct mathematical instruction tuning on DeepSeekMath-Base utilizing chain-of-thought \citep{cot}, program-of-thought \citep{pot,pal}, and tool-integrated reasoning \citep{tora} datasets. The tuned model, DeepSeekMath-Instruct 7B, surpasses all other 7B models and demonstrates performance on par with instruction-tuned open-source models of 70B scale.

In addition, we present Group Relative Policy Optimization (GRPO), a new variant of reinforcement learning (RL) inspired by Proximal Policy Optimization (PPO) \citep{schulman2017proximal}. Unlike PPO, GRPO eliminates the need for a critic model and instead utilizes group scores to estimate the baseline, leading to a marked reduction in the computational overhead during training. Applying GRPO exclusively on a portion of English instruction tuning data yields notable gains over the already robust DeepSeekMath-Instruct baseline; in the reinforcement learning stage, it delivers clear improvements on both familiar datasets (GSM8K: 82.9\% $\rightarrow$ 88.2\%, MATH: 46.8\% $\rightarrow$ 51.7\%) and unfamiliar mathematical benchmarks (e.g., CMATH: 84.6\% $\rightarrow$ 88.8\%). We further contribute a unified framework that contextualizes approaches such as Rejection Sampling Fine-Tuning (RFT) \citep{yuan2023scaling}, Direct Preference Optimization (DPO) \citep{dpo}, PPO, and GRPO. Through this integrative lens, we observe that each method may be categorized as a form of either direct or simplified RL. To further elucidate the key components underpinning this framework, we perform comprehensive experiments examining variables such as online versus offline training, outcome-oriented versus process-oriented supervision, and single-step versus iterative RL. Finally, we analyze the reasons underlying the efficacy of our RL approach in enhancing instruction-tuned model performance, and propose promising avenues for advancing RL methodologies within this unified scheme.

\subsection{Contributions}
We present scalable approaches to math pre-training and conduct an in-depth investigation and assessment of reinforcement learning methods.

\textbf{Math Pre-Training at Scale}

\begin{itemize}[topsep=0pt]
    \item This study demonstrates that the openly available Common Crawl dataset holds substantial resources relevant to mathematical research. Through the application of a carefully structured selection process, we assemble the DeepSeekMath Corpus—a high-caliber collection comprising 120B tokens drawn from mathematically relevant web pages. This corpus surpasses the size of Minerva's math web data by nearly sevenfold \citep{minerva}, and exceeds the scale of OpenWebMath \citep{openwebmath}, released more recently, by a factor of nine.

\item The DeepSeekMath-Base 7B pre-trained model demonstrates performance on par with Minerva 540B \citep{minerva}, suggesting that sheer parameter count does not solely determine mathematical reasoning proficiency. Robust results can also be obtained from more compact models when they are pre-trained using high-quality datasets.

\item We present results from our experiments focused on math training. When models undergo code training before math training, their performance in solving mathematical tasks is enhanced, regardless of whether tools are involved. This outcome provides partial insight into the enduring inquiry: \textit{does code training enhance reasoning skills?} Our evidence suggests that it does—specifically in the domain of mathematical reasoning.

\item While the use of arXiv papers for training is prevalent, particularly among various mathematics-focused studies, it does not yield significant gains across any of the mathematical benchmarks utilized in this work.
\end{itemize}

\textbf{Exploration and Analysis of Reinforcement Learning}

\begin{itemize}[topsep=0pt]
    \item We present Group Relative Policy Optimization (GRPO), a reinforcement learning method designed for both efficiency and efficacy. GRPO eliminates the need for a critic model by estimating baselines through group scores, thus greatly decreasing the computational demands when compared to Proximal Policy Optimization (PPO).
    \item Our findings reveal that integrating GRPO with instruction-tuning data leads to notable improvements in the performance of DeepSeekMath-Instruct, our instruction-aligned model. Additionally, the approach also fosters improvements on out-of-domain generalization during the reinforcement learning phase.
    \item We establish a cohesive framework for understanding a range of methods, including RFT, DPO, PPO, and GRPO. Furthermore, we carry out comprehensive experimental analyses encompassing comparisons such as online versus offline learning, supervision at the outcome versus process levels, as well as single-turn versus iterative reinforcement learning, to thoroughly examine the fundamental components of this framework.
    \item Utilizing the proposed unified paradigm, we analyze the underlying mechanisms driving the success of reinforcement learning and outline several promising avenues for advancing the reinforcement learning of large language models (LLMs).
\end{itemize}

\subsection{Summary of Evaluations and Metrics}

\begin{itemize}[topsep=0pt]
    \item \textbf{English and Chinese Mathematical Reasoning}: We systematically evaluate our models using a diverse suite of benchmarks in both English and Chinese, spanning mathematical tasks from elementary through collegiate levels. English evaluations utilize datasets such as GSM8K \citep{gsm8k}, MATH \citep{MATH}, SAT \citep{llemma}, OCW Courses \citep{minerva}, and MMLU-STEM \citep{mmlu}. For Chinese, we employ benchmarks like MGSM-zh \citep{mgsm}, CMATH \citep{wei2023cmath}, Gaokao-MathCloze \citep{agieval}, and Gaokao-MathQA \citep{agieval}. Our analysis covers both the generation of independent written solutions (without external tools) and problem-solving through the application of Python.

On a range of English evaluation tasks, DeepSeekMath-Base demonstrates performance on par with the proprietary Minerva 540B model \citep{minerva}, and consistently outperforms all available open-source base models—including Mistral 7B \citep{mistral} and Llemma-34B \citep{llemma}—irrespective of their engagement in math-focused pre-training, frequently by a wide margin. Remarkably, DeepSeekMath-Base achieves even better results on Chinese assessment tasks. This may be attributed to our divergence from prior methodologies \citep{minerva,llemma}, as we expand our pre-training corpus beyond English, incorporating high-quality mathematical data in other languages. Leveraging subsequent instruction tuning and reinforcement learning in math, DeepSeekMath-Instruct and DeepSeekMath-RL show substantial improvements, achieving—for the first time among open-source models—accuracy rates exceeding 50\% on the highly challenging, competition-grade MATH dataset.

\item \textbf{Formal Mathematics}: DeepSeekMath-Base is assessed on its informal-to-formal theorem proving ability, following the protocol of \citep{dsp_proof}, utilizing the miniF2F benchmark \citep{minif2f} with Isabelle \citep{isabelle} selected as the proof assistant. The model exhibits robust performance in few-shot autoformalization scenarios.
\item \textbf{Natural Language Understanding, Reasoning, and Code}: For a thorough evaluation of the model's capacities in language understanding, reasoning, and programming, DeepSeekMath-Base is tested on the Massive Multitask Language Understanding (MMLU) suite \citep{mmlu}, which spans 57 multiple-choice categories, and BIG-Bench Hard (BBH) \citep{bbh}, comprised of 23 advanced tasks centered primarily on complex, multi-step reasoning. In addition, the model is benchmarked on HumanEval \citep{codex} and MBPP \citep{mbpp}, both standard datasets for assessing coding proficiency in language models. Mathematical pre-training enhances both linguistic comprehension and reasoning accuracy.

\section{Math Pre-Training}

\subsection{Data Collection and Decontamination} This section describes the methodology used to develop the DeepSeekMath Corpus utilizing data from Common Crawl. Illustrated in Figure \ref{fig:data_collect}, the process involves an iterative pipeline designed to efficiently assemble an extensive mathematical dataset beginning with an initial seed corpus (such as a curated, high-quality set of math-focused data). Importantly, this framework is generalizable and can be extended to additional fields, such as programming.

Our methodology begins with OpenWebMath \citep{openwebmath}, a curated dataset comprising high-quality mathematical web resources, serving as the foundation for our seed corpus. Leveraging this corpus, we utilize the fastText model \citep{joulin2016fasttext} to retrieve additional mathematical web pages with characteristics similar to those in OpenWebMath. For training, we randomly extract 500,000 examples from the seed corpus as positive samples, alongside 500,000 web pages drawn from Common Crawl as negatives. Training is carried out using an open-source implementation\footnote{\url{https://fasttext.cc}}, where we set the embedding size to 256, adopt a learning rate of 0.1, utilize word n-grams up to length 3, require a minimum word frequency of 3, and run training for 3 epochs. In order to minimize redundancies in the original Common Crawl data, we apply both URL-based deduplication and techniques for detecting near-duplicate content, reducing the total to 40 billion HTML pages. Using the trained fastText model, we identify candidate mathematical web pages within this deduplicated dataset. We then employ the fastText prediction scores to sort the retrieved pages and retain only the highest scoring subset, thereby filtering out material of lower relevance or quality. The final selection of data volume is determined through pre-training runs that examine the performance on datasets containing the top 40B, 80B, 120B, and 160B tokens. For the initial pass, we opt to retain the leading 40B tokens.

Following the initial round of data gathering, a significant number of mathematical web pages remain uncollected. This gap primarily arises due to the limited diversity in the set of positive examples used to train the fastText model. To address this issue, we broaden the range of mathematical sources included in the seed corpus, thereby enhancing the training of the fastText model. Our process begins by partitioning the entire Common Crawl into non-overlapping domains, each defined by a unique base URL. For every domain, we determine the proportion of web pages that were gathered during the first round. If a domain has more than 10\% of its web pages collected, it is categorized as math-oriented (for instance, \url{mathoverflow.net}). 

Next, we carry out manual annotation of URLs within these math-related domains that correspond to mathematical content (such as \url{mathoverflow.net/questions}). Any web pages associated with these mathematical URLs that have yet to be collected are incorporated into the seed corpus. This strategy yields an expanded set of positive examples, equipping the fastText model to retrieve a broader scope of mathematical web content in the subsequent data collection cycle. Over the course of four iterations, this method yields a total of 35.5 million mathematical web pages, amounting to 120 billion tokens. By the conclusion of the fourth round, we observe that almost 98\% of the mathematical data identified had already been acquired by the third iteration, leading us to discontinue further data collection.

In order to prevent contamination of our benchmarks, we adopt the methodology outlined in \cite{deepseek-coder}, filtering out web content that includes questions or answers drawn from prominent English mathematical benchmarks such as GSM8K \citep{gsm8k} and MATH \citep{MATH}, as well as from Chinese benchmarks like CMATH \citep{wei2023cmath} and AGIEval \citep{agieval}. The specific approach we employ is to exclude any portion of text containing a 10-gram sequence that precisely matches any substring from these evaluation benchmarks. Additionally, for those benchmark excerpts comprising fewer than 10 but at least 3 grams, we apply strict exact-match filtering to eliminate affected web pages from our mathematical training dataset.

\subsection{Validating the Quality of the DeepSeekMath~Corpus}

We run pre-training experiments to investigate how the DeepSeekMath~Corpus is compared with the recently released math-training corpora:

\begin{itemize}[topsep=0pt]
    \item \textbf{MathPile} \citep{mathpile}: a collection merging data from diverse sources (8.9B tokens), including textbooks, Wikipedia, ProofWiki, CommonCrawl, StackExchange, and arXiv—with arXiv constituting the predominant portion (over 85\%) of the dataset;
    \item \textbf{OpenWebMath} \citep{openwebmath}: mathematical content filtered out of CommonCrawl and comprising a total of 13.6B tokens;
    \item \textbf{Proof-Pile-2} \citep{llemma}: a mathematical dataset that integrates OpenWebMath, AlgebraicStack (10.3B tokens from mathematical code), and papers from arXiv (28.0B tokens). In our usage of Proof-Pile-2, we adhere to the arXiv:Web:Code ratio of 2:4:1 as established by \cite{llemma}.
\end{itemize}

\subsubsection{Training Setting}
\label{sec:quality-policy}
We perform mathematical instruction tuning on a general-purpose language model comprising 1.3B parameters, architecturally consistent with the DeepSeek LLMs \citep{deepseek-llm}—this model is referenced as DeepSeek-LLM 1.3B. Distinct instances of the model are fine-tuned individually on each specialized mathematical dataset, for a total of 150B tokens per instance. All training processes leverage the resource-efficient HAI-LLM \citep{haillm} system. Consistent with the established procedures in the DeepSeek LLMs literature, training utilizes the AdamW optimizer \citep{adamW}, configured with $\beta_1=0.9$, $\beta_2=0.95$, and $\mathrm{weight\_decay}=0.1$. The learning rate follows a multi-stage decay: after an initial linear warmup phase of 2,000 steps to the maximal value, it is reduced to 31.6\% at the 80\% progress mark, and subsequently to 10.0\% at 90\% of the total training duration. The peak learning rate is fixed at $5.3 \times 10^{-4}$, while each mini-batch contains 4M tokens within a 4K token context window.

\subsubsection{Evaluation Results}

\textbf{The DeepSeekMath~Corpus is of high quality, covers multilingual mathematical content, and is the largest in size.}

\begin{itemize}[topsep=0pt]
    \item \textbf{High-quality}:
    To assess downstream capabilities, we conduct few-shot chain-of-thought prompting \cite{cot} across 8 mathematical benchmarks. Table \ref{tab:corpora_comparison} illustrates a distinct performance advantage for models trained using the DeepSeekMath~Corpus. Moreover, Figure \ref{fig:corpora_comparisons} reveals that the DeepSeekMath Corpus model surpasses Proof-Pile-2’s results at the 50B token point (corresponding to one complete Proof-Pile-2 epoch), which suggests the DeepSeekMath Corpus offers superior average data quality.
    \item \textbf{Multilingual}:
    The DeepSeekMath~Corpus contains multilingual content, with English and Chinese being the most prevalent languages. Training with the DeepSeekMath~Corpus, as indicated in Table \ref{tab:corpora_comparison}, yields gains in mathematical reasoning for both English and Chinese tasks. By contrast, other existing mathematical datasets are predominantly English-oriented, which restricts their efficacy and can even negatively affect performance on Chinese mathematical problems.
    \item \textbf{Large-scale}:
    In terms of data volume, the DeepSeekMath~Corpus significantly exceeds existing mathematical corpora in size. Figure \ref{fig:corpora_comparisons} shows that DeepSeek-LLM 1.3B achieves a sharper and more sustained trajectory of improvement when trained on DeepSeekMath~Corpus. In comparison, baseline datasets are markedly smaller and have been cycled through the training process multiple times, resulting in the model’s learning stagnating after a short period.
\end{itemize}

\subsection{Training and Evaluating DeepSeekMath-Base 7B}

In this section, we present DeepSeekMath-Base 7B, a foundational model demonstrating robust mathematical reasoning capabilities. The model's initial weights are derived from DeepSeek-Coder-Base-v1.5 7B \citep{deepseek-coder}, and its training encompasses 500B tokens. Data utilized in training are distributed as follows: 56\% originates from the DeepSeekMath Corpus, 4\% from AlgebraicStack, 10\% from arXiv papers, 20\% from Github source code, and the remaining 10\% consists of natural language inputs in English and Chinese from Common Crawl. With respect to training procedures, we largely follow the protocols outlined in Section \ref{sec:quality-policy}, though we adjust the maximum learning rate to 4.2e-4 and set the batch size to 10 million tokens.

This work presents an in-depth evaluation of the mathematical proficiency of DeepSeekMath-Base 7B, emphasizing its capacity to generate independent mathematical solutions without assistance from external resources, tackle mathematical tasks utilizing auxiliary tools, and perform rigorous formal theorem proving. In addition to these mathematical assessments, we extend our analysis to encompass a broader overview of the base model’s general performance, examining its abilities in natural language understanding, logical reasoning, and programming.

\paragraph{Mathematical Problem Solving with Step-by-Step Reasoning}

We assess the capabilities of DeepSeekMath-Base in addressing mathematical problems using few-shot chain-of-thought prompting \citep{cot} over eight distinct English and Chinese benchmarks. These benchmarks span areas such as quantitative reasoning—including GSM8K \citep{gsm8k}, MATH \citep{MATH}, and CMATH \citep{wei2023cmath}—as well as multiple-choice tasks, like those found in MMLU-STEM \citep{mmlu} and Gaokao-MathQA \citep{agieval}. Collectively, these datasets represent a broad spectrum of mathematical disciplines, ranging from elementary mathematics to problems of collegiate difficulty.

As presented in Table \ref{tab:base_math_cot}, DeepSeekMath-Base 7B demonstrates superior results across all eight evaluation benchmarks when compared to other open-source base models, such as the popular Mistral 7B \citep{mistral} and the recently introduced Llemma 34B \citep{llemma}, which has been further trained on Proof-Pile-2 \citep{llemma} for mathematics. Remarkably, on the highly challenging MATH dataset, DeepSeekMath-Base exceeds the performance of all available open-source base models by a margin greater than 10\% in absolute terms. Furthermore, it achieves better results than Minerva 540B \citep{minerva}, a much larger proprietary model—77 times the size—that extends PaLM \citep{palm} with additional training on mathematical materials.

\paragraph{Mathematical Problem Solving with Tool Use} We assess the ability of models to perform mathematical reasoning supported by external tools on GSM8K and MATH benchmarks, employing the few-shot program-of-thought prompting methodology described by \citet{pot,pal}. In these experiments, models are asked to generate Python scripts to solve given problems, leveraging libraries like \textit{math} and \textit{sympy} for handling complex calculations. The answer is then taken from the output of the executed program. According to the results presented in Table \ref{tab:base_math_pot_proof}, DeepSeekMath-Base 7B surpasses the previous leading model, Llemma 34B.

\paragraph{Formal Mathematics}

The automation of formal proofs plays a crucial role in verifying the correctness and robustness of mathematical arguments while also streamlining the proof development process, a topic that has garnered considerable interest recently. In this work, we assess the capability of DeepSeekMath-Base 7B on the informal-to-formal proof generation task described by \citep{dsp_proof}, wherein the objective is to produce a formal proof given an informal problem statement, its formal analogue, and an informal justification. Our evaluation utilizes the miniF2F benchmark \citep{minif2f}, which targets formalizations of mathematical problems at the Olympiad level. We employ few-shot prompting to produce Isabelle formal proofs for each selected problem. Consistent with the methodology in \cite{dsp_proof}, proof sketches are generated by the language model, and Sledgehammer \citep{sledgehammer}, a standard automated proving tool, is invoked to complete any gaps. The outcomes, as presented in Table \ref{tab:base_math_pot_proof}, indicate that DeepSeekMath-Base 7B achieves notable results in the autoformalization of proofs.

\paragraph{Natural Language Understanding, Reasoning, and Code}

We assess the performance of models in natural language understanding using MMLU \citep{mmlu}, reasoning ability via BBH \citep{bbh}, and programming skills with HumanEval \citep{codex} and MBPP \citep{mbpp}. Table \ref{tab:understanding_reasoning_code} indicates that DeepSeekMath-Base 7B achieves marked improvements on MMLU and BBH compared to its predecessor, DeepSeek-Coder-Base-v1.5 \citep{deepseek-coder}, highlighting the advantageous influence of mathematical training for tasks involving language comprehension and reasoning. Furthermore, by continually incorporating code tokens into training, DeepSeekMath-Base 7B is able to retain the coding performance demonstrated by DeepSeek-Coder-Base-v1.5 on both coding evaluations. In summary, DeepSeekMath-Base 7B provides a substantial performance gain over the generalist Mistral 7B model \citep{mistral} across the three evaluated reasoning and coding benchmarks.

\section{Supervised Fine-Tuning}

\subsection{SFT Data Curation}

We develop a comprehensive mathematical instruction-tuning dataset that encompasses both English and Chinese questions drawn from a variety of mathematical domains and exhibiting a range of difficulty levels. Each problem is matched with a corresponding solution presented in formats such as chain-of-thought (CoT) \citep{cot}, program-of-thought (PoT) \citep{pot,pal}, as well as tool-integrated reasoning \citep{tora}. In total, the dataset includes 776K training instances.

\begin{itemize}[topsep=0pt]
    \item \textbf{English mathematical datasets}: We provide tool-augmented solution annotations for problems from GSM8K and MATH, and incorporate selected samples from MathInstruct \citep{MathInstruct}, as well as examples from the Lila-OOD training set \citep{lila} that feature solutions utilizing CoT or PoT approaches. Our compilation of English data encompasses a broad array of mathematical disciplines including algebra, probability, number theory, calculus, and geometry.
    \item \textbf{Chinese mathematical datasets}: Our collection consists of Chinese K-12 mathematics questions covering 76 distinct sub-domains such as linear equations, each annotated with detailed solutions in both CoT and tool-assisted reasoning styles.
\end{itemize}

\subsection{Training and Evaluating DeepSeekMath-Instruct 7B}

This section presents DeepSeekMath-Instruct 7B, developed via mathematical instruction tuning from the DeepSeekMath-Base model. For training, samples are randomly combined together up to a maximum context length of 4K tokens. The model is trained for 500 iterations, utilizing a batch size of 256 and a fixed learning rate set at 5e-5.

We evaluate models' mathematical performance both without and with tool use, on 4 quantitative reasoning benchmarks in English and Chinese. We benchmark our model against the leading models of the time:

\begin{itemize}[topsep=0pt]
    \item \textbf{Closed-source models} encompass: (1) the GPT series, notably GPT-4 \citep{gpt4} and the GPT-4 Code Interpreter~\footnote{\url{https://openai.com/blog/chatgpt-plugins\#\#code-interpreter}}, recognized for their advanced capabilities, (2) Gemini Ultra and Gemini Pro \citep{gemini}, (3) Inflection-2 \citep{inflection-2}, (4) Grok-1~\footnote{\url{https://x.ai/model-card}}, and a selection of recently introduced models from China, including (5) Baichuan-3~\footnote{\url{https://www.baichuan-ai.com}}, and (6) the most recent GLM-4~\footnote{\url{https://open.bigmodel.cn/dev/api\#glm-4}}

    from the GLM family \citep{glm}.

These models are intended for broad application, with most having been subjected to various alignment strategies.

\item \textbf{Open-source models} encompass: general-purpose models such as (1) DeepSeek-LLM-Chat 67B \citep{deepseek-llm}, (2) Qwen 72B \citep{qwen}, (3) SeaLLM-v2 7B \citep{seallm}, and (4) ChatGLM3 6B \citep{chatglm3}. Additionally, there are models tailored for improved mathematical capabilities, including (5) InternLM2-Math 20B~\footnote{\url{https://github.com/InternLM/InternLM-Math}}, which is based on InternLM2 and receives math-specific training prior to instruction fine-tuning, (6) Math-Shepherd-Mistral 7B, where PPO training \citep{schulman2017proximal} is applied to Mistral 7B \citep{mistral} using a reward model guided by process supervision, (7) the WizardMath series \citep{wizardmath}, which enhances mathematical reasoning in Mistral 7B and Llama-2 70B \citep{llama2} by leveraging evolve-instruct (an AI-generated instruction tuning strategy) and PPO-based learning, with training samples largely drawn from GSM8K and MATH, (8) MetaMath 70B \citep{metamath}, which involves Llama-2 70B being fine-tuned on enriched GSM8K and MATH datasets, (9) ToRA 34B \cite{tora}, which modifies CodeLlama 34B through additional training to perform tool-augmented mathematical reasoning, and (10) MAmmoTH 70B \citep{MathInstruct}, where Llama-2 70B is instruction-tuned using the MathInstruct dataset.

Referencing Table \ref{tab:sft_rl_math}, when evaluated without the assistance of external tools, DeepSeekMath-Instruct 7B achieves impressive results in step-by-step mathematical reasoning. Remarkably, on the MATH dataset, which targets competition-level problems, our model outperforms every available open-source alternative as well as most proprietary systems—including Inflection-2 and Gemini Pro—by a margin of no less than 9\% in absolute terms. This performance gap holds even when compared with much larger models (such as Qwen 72B) and those optimized explicitly via math-oriented reinforcement learning techniques (for instance, WizardMath-v1.1 7B). Although DeepSeekMath-Instruct demonstrates competitiveness with leading Chinese commercial models like GLM-4 and Baichuan-3 on this benchmark, it is still outperformed by models such as GPT-4 and Gemini Ultra.

In the evaluation scenario permitting the combination of natural language reasoning with programmatic tool utilization for solving problems, DeepSeekMath-Instruct 7B achieves nearly 60\% accuracy on MATH, outperforming all previously released open-source models. Across additional benchmarks, our model demonstrates performance on par with DeepSeek-LLM-Chat 67B, the preceding state-of-the-art system, despite the latter possessing ten times more parameters.
\section{Reinforcement Learning}

\subsection{Group Relative Policy Optimization} Reinforcement learning (RL) has demonstrated considerable success in enhancing the mathematical reasoning capabilities of LLMs following the Supervised Fine-Tuning (SFT) phase \citep{wang2023math,wizardmath}. Here, we present Group Relative Policy Optimization (GRPO), a novel RL approach designed to be both efficient and effective.

\subsubsection{From PPO to GRPO} Proximal Policy Optimization (PPO) \citep{schulman2017proximal} is an actor-critic RL algorithm that is widely used in the RL fine-tuning stage of LLMs \citep{ouyang2022training}. In particular, it optimizes LLMs by maximizing the following surrogate objective:

\begin{equation}
\footnotesize
    \mathcal{J}_{PPO}(\theta) = \mathbb{E}_{q \sim P(Q),\ o \sim \pi_{\theta_{old}}(O|q)} \left[ \frac{1}{|o|} \sum_{t=1}^{|o|} \min \left( \frac{\pi_\theta(o_{t} | q, o_{<t})}{\pi_{\theta_{old}}(o_{t} | q, o_{<t})} A_{t}, \text{clip} \left( \frac{\pi_\theta(o_{t} | q, o_{<t})}{\pi_{\theta_{old}}(o_{t} | q, o_{<t})}, 1 - \varepsilon, 1 + \varepsilon \right)  A_{t} \right) \right],
\end{equation}

Here, $\pi_{\theta}$ refers to the current policy model, while $\pi_{\theta_{old}}$ denotes the previous version of the policy. The variables $q$ and $o$ correspond to questions drawn from the question dataset and outputs sampled from the old policy $\pi_{\theta_{old}}$, respectively. The parameter $\varepsilon$ is a hyper-parameter associated with clipping, introduced in PPO to promote stable training dynamics. The advantage $A_t$ is determined using Generalized Advantage Estimation (GAE) \citep{gae}, calculated from the rewards $\{r_{\ge t}\}$ in conjunction with a learned value function $V_{\psi}$. As a result, PPO necessitates training a value function simultaneously with the policy. To prevent excessive reward model optimization, it is standard practice to incorporate a KL divergence penalty from a reference model into the reward at every token, applied on a per-token basis \citep{ouyang2022training}, such that

\begin{equation}
    r_{t} = r_\varphi(q, o_{\le t}) - \beta \log\frac{\pi_{\theta}(o_{t}|q, o_{<t})}{\pi_{ref}(o_{t}|q, o_{<t})},
\label{eq:PPO-reward}
\end{equation}

Here, $r_\varphi$ denotes the reward model, $\pi_{ref}$ represents the reference model—typically the original SFT model—and $\beta$ serves as the weight for the KL divergence penalty.

As the value function employed in PPO is typically another model of comparable size as the policy model, it brings a substantial memory and computational burden. Additionally, during RL training, the value function is treated as a baseline in the calculation of the advantage for variance reduction. While in the LLM context, usually only the last token is assigned a reward score by the reward model, which may complicate the training of a value function that is accurate at each token. To address this, as shown in Figure \ref{fig:grpo}, we propose Group Relative Policy Optimization (GRPO), which obviates the need for additional value function approximation as in PPO, and instead uses the average reward of multiple sampled outputs, produced in response to the same question, as the baseline. More specifically, for each question $q$, GRPO samples a group of outputs $\{o_1, o_2, \cdots, o_G\}$  from the old policy  $\pi_{\theta_{old}}$  and then optimizes the policy model by maximizing the following objective:

\begin{equation}
\footnotesize
\begin{split}
    \mathcal{J}_{GRPO}(\theta) &= \mathbb{E}{[q \sim P(Q), \{o_i\}_{i=1}^G \sim \pi_{\theta_{old}}(O|q)]}  \\
    & \frac{1}{G}\sum_{i=1}^G\frac{1}{|o_i|} \sum_{t=1}^{|o_i|} \bigg\{ \min \left[ \frac{\pi_\theta(o_{i,t} | q, o_{i,<t})}{\pi_{\theta_{old}}(o_{i,t} | q, o_{i,<t})} \hat{A}_{i,t}, \text{clip} \left( \frac{\pi_\theta(o_{i,t} | q, o_{i,<t})}{\pi_{\theta_{old}}(o_{i,t} | q, o_{i,<t})}, 1 - \varepsilon, 1 + \varepsilon \right)  \hat{A}_{i,t} \right] \\
    & \hspace{6cm} - \beta \mathbb{D}_{KL}\left[\pi_{\theta} || \pi_{ref}\right]\bigg\} ,
\end{split}
\label{eq:GRPO-obj}
\end{equation}

where $\varepsilon$ and $\beta$ are hyper-parameters, and $\hat{A}_{i,t}$  is the advantage calculated based on relative rewards of the outputs inside each group only, which will be detailed in the following subsections. The group relative way that GRPO leverages to calculate the advantages, aligns well with the comparative nature of rewards models,  as reward models are typically trained on datasets of comparisons between outputs on the same question. Also note that, instead of adding KL penalty in the reward, GRPO regularizes by directly adding the KL divergence between the trained policy and the reference policy to the loss, avoiding complicating the calculation of  $\hat{A}_{i,t}$. And different from the KL penalty term used in (\ref{eq:PPO-reward}), we estimate the KL divergence with the following unbiased estimator \citep{kl_approx}:

\begin{equation}
\small
    \mathbb{D}_{KL}\left[\pi_{\theta} || \pi_{ref}\right] = \frac{\pi_{ref}(o_{i,t}|q,o_{i,<t})}{\pi_{\theta}(o_{i,t}|q,o_{i,<t})}- \log\frac{\pi_{ref}(o_{i,t}|q,o_{i,<t})}{\pi_{\theta}(o_{i,t}|q,o_{i,<t})} - 1,
\end{equation}

which is guaranteed to be positive.

\begin{algorithm}[t]
  \small
  \caption{Iterative Group Relative Policy Optimization}
  \textbf{Input} starting policy model $\pi_{\theta_{\text{init}}}$; reward models $r_\varphi$; collection of task prompts $\mathcal{D}$; 
  hyperparameters $\varepsilon$, $\beta$, $\mu$
  \begin{algorithmic}[1]
    \State Set policy model $\pi_\theta \leftarrow \pi_{\theta_{\text{init}}}$
    \For{iteration = 1, \dots, I}
      \State Assign reference model $\pi_{ref} \leftarrow \pi_{\theta}$
      \For{step = 1, \dots, M}
      \State Draw mini-batch $\mathcal{D}_b$ from $\mathcal{D}$
      \State Replace previous policy model $\pi_{\theta_{old}} \leftarrow \pi_{\theta}$
      \State For each query $q \in \mathcal{D}_b$, sample $G$ outputs $\{o_i\}_{i=1}^G \sim \pi_{\theta_{old}} (\cdot \mid q)$
      \State For every sampled output $o_i$, calculate corresponding rewards $\{r_i\}_{i=1}^G$ using $r_{\varphi}$
      \State Compute $\hat{A}_{i,t}$ for the $t$-th token of each $o_i$ using group-relative advantage calculation
      \For{each GRPO update = 1, \dots, $\mu$}
        \State Modify policy model $\pi_{\theta}$ by optimizing the GRPO loss (see Equation \ref{eq:GC-GRPO})
      \EndFor
    \EndFor
    \State Continually adapt $r_\varphi$ using a replay-based training strategy. 
    \EndFor 
  \end{algorithmic}
  \textbf{Output} $\pi_\theta$
  \label{alg:iter-grpo}
\end{algorithm}

\subsubsection{Outcome Supervision RL with GRPO}  
Specifically, given a question $q$, a collection of responses $\{o_1, o_2, \cdots, o_G\}$ are generated using the previous iteration of the policy, denoted as $\pi_{\theta_{old}}$. These outputs are each evaluated by a reward model, which assigns a set of $G$ reward values $\mathbf{r} = \{r_1, r_2, \cdots, r_G\}$ accordingly. To standardize the rewards, each is adjusted by first subtracting the mean reward of the group and then dividing by the group's standard deviation. Under outcome supervision, this normalized reward is provided only at the conclusion of each response $o_i$, and the computed advantages $\hat{A}_{i, t}$ for all tokens in $o_i$ are set equal to the normalized reward. That is, $\hat{A}_{i, t} = \widetilde{r}_i = \frac{r_i- {\rm mean}(\mathbf{r})}{{\rm std}(\mathbf{r})}$. Policy optimization then proceeds by maximizing the objective specified in equation (\ref{eq:GRPO-obj}).

\subsubsection{Process Supervision RL with GRPO}

Outcome supervision typically only assigns a reward at the conclusion of an output, which can be inadequate and inefficient for directing the policy in intricate mathematical problems. In line with the methodology in \cite{wang2023math}, our approach also considers process supervision, in which a reward is given following each step of reasoning. Specifically, for a given question $q$ and $G$ sampled outputs $\{o_1, o_2, \cdots, o_G\}$, a process reward model evaluates each individual step within the outputs, producing a set of rewards structured as follows: $\mathbf{R} = \{ \{r_1^{index(1)},\cdots,r_1^{index(K_1)}\}, \cdots,  \{r_G^{index(1)},\cdots,r_G^{index(K_G)}\} \}$. Here, $index(j)$ corresponds to the ending token position of the $j$-th step, and $K_i$ signifies the number of steps present in the $i$-th output. These reward values are standardized by subtracting the mean and dividing by the standard deviation, yielding $\widetilde{r}_i^{index(j)} = \frac{r_i^{index(j)} - {\rm mean(\mathbf{R})}}{{\rm std(\mathbf{R})}}$.

Thereafter, in process supervision, the advantage associated with each token is computed as the sum of the normalized rewards for all subsequent steps, formalized by $\hat{A}_{i, t} = \sum_{index(j) \ge t} \widetilde{r}_i^{index(j)}$. The policy is then trained to maximize the objective specified in equation (\ref{eq:GRPO-obj}).

\subsubsection{Iterative RL with GRPO} During the course of reinforcement learning, the existing reward model may eventually fall short in effectively guiding the updated policy model. To address this, we investigate iterative RL with GRPO. As detailed in Algorithm \ref{alg:iter-grpo}, this approach involves regularly constructing new datasets for the reward model using samples obtained from the policy model. The reward model is then updated iteratively, employing a replay strategy that retains 10\% of previous data to maintain historical context. Subsequently, the reference model is aligned with the policy model, which is then further optimized using the latest version of the reward model.

\subsection{Training and Evaluating DeepSeekMath-RL}

Our reinforcement learning (RL) experiments utilize DeepSeekMath-Instruct 7B as the underlying model. The RL training data comprises chain-of-thought-style items drawn from GSM8K and MATH within the SFT corpus, totaling approximately 144K questions. Other SFT-derived problems are purposefully omitted to specifically analyze how RL influences benchmarks that have no data represented during RL training. The reward model’s training dataset is assembled in alignment with procedures outlined by \citep{wang2023math}. The initial reward model is trained on DeepSeekMath-Base 7B, employing a learning rate of 2e-5. During GRPO training, the policy model’s learning rate is fixed at 1e-6, and a KL penalty coefficient of 0.04 is used. We generate $64$ samples per prompt, employ a sequence length cap of 1024 tokens, and set the training batch size at 1024. The policy model receives just one update after each exploration step. Evaluation of DeepSeekMath-RL 7B is carried out on the same test sets as DeepSeekMath-Instruct 7B. Notably, for DeepSeekMath-RL 7B, GSM8K and MATH in chain-of-thought format are categorized as in-domain, whereas all other evaluation sets are considered out-of-domain.

Table \ref{tab:sft_rl_math} presents the comparative results of both open- and closed-source models engaging in chain-of-thought and tool-integrated reasoning across English and Chinese evaluation tasks. The findings are as follows: (1) When applying chain-of-thought reasoning, DeepSeekMath-RL 7B achieves accuracies of 88.2\% on GSM8K and 51.7\% on MATH, outperforming all open-source competitors within the 7B to 70B scale, as well as most closed-source systems. (2) Notably, DeepSeekMath-RL 7B is exclusively fine-tuned on chain-of-thought-format instructional data from GSM8K and MATH, building upon DeepSeekMath-Instruct 7B. Even with such a limited training dataset, DeepSeekMath-RL 7B consistently exceeds the performance of DeepSeekMath-Instruct 7B across every assessment metric, indicating the substantial benefit conferred by reinforcement learning.

\section{Discussion}
This section presents an overview of our observations derived from both pre-training and reinforcement learning experiments.

\subsection{Lessons Learnt in Pre-Training}

We begin by presenting our observations during the pre-training phase. Except where explicitly mentioned, all training configurations correspond to those described in Section \ref{sec:quality-policy}. Additionally, it should be emphasized that references to the DeepSeekMath Corpus throughout this section pertain to a dataset comprising 89 billion tokens, collected during the second round of data gathering.

\subsubsection{Code Training Benefits Mathematical Reasoning}

A widely discussed but as yet unproven hypothesis posits that code training enhances reasoning abilities. In this work, we seek to partially address this claim within the context of mathematics, finding that code training bolsters models' mathematical reasoning skills, irrespective of whether external tools are employed.

To study how code training affects mathematical reasoning, we experimented with the following two-stage training and one-stage training settings:

\textbf{Two-Stage Training}

\begin{itemize}[topsep=0pt]
    \item \textbf{Code Training for 400B Tokens $\rightarrow$ Math Training for 150B Tokens}: We initiate training of DeepSeek-LLM 1.3B using 400B tokens drawn from code data, subsequently transitioning to an additional 150B tokens focused on mathematics;
    \item \textbf{General Training for 400B Tokens $\rightarrow$ Math Training for 150B Tokens}: For comparison, a parallel experiment is conducted where the initial 400B tokens are from a broad general corpus (curated by DeepSeek-AI) instead of code, followed by identical mathematics training. This setup aims to discern whether code-based pretraining affords benefits to mathematical reasoning relative to general-domain pretraining.
\end{itemize}

\textbf{One-Stage Training}

\begin{itemize}[topsep=0pt]
    \item \textbf{Math-Focused Training on 150B Tokens}: DeepSeek-LLM 1.3B is optimized using 150 billion math-specific tokens;
    \item \textbf{Joint Training Utilizing 400B Code Tokens and 150B Math Tokens}: We observe that sequentially training on math tokens after code tokens can negatively impact code performance. Therefore, we explore whether simultaneously integrating code and math tokens in a unified training phase enhances mathematical reasoning while also reducing the risk of catastrophic forgetting.
\end{itemize}

\paragraph{Results}
The downstream outcomes across various training configurations are presented in Table \ref{tab:code-math} and Table \ref{tab:code-math-general-eval}.

Integrating code training into the curriculum supports program-aided mathematical reasoning across both two-stage and one-stage training paradigms. As illustrated in Table \ref{tab:code-math}, implementing code training in the two-stage approach already provides a notable boost to Python-based GSM8K and MATH problem-solving capabilities. Introducing math training during the subsequent stage contributes additional gains. Notably, within a one-stage training framework, combining code and math tokens prevents the catastrophic forgetting commonly observed with two-stage training, while also enhancing coding proficiency (Table \ref{tab:code-math-general-eval}) and supporting program-based mathematical reasoning (Table \ref{tab:code-math}).

Engaging in code training yields gains in mathematical reasoning abilities even in the absence of tools. Within a two-stage training framework, introducing code training at the outset provides moderate improvements. This approach further facilitates more effective math training in the following phase, culminating in optimal performance. Conversely, merging code tokens with math tokens during a single-stage training regimen diminishes mathematical reasoning skills when tools are not employed. A possible explanation is that, given the relatively small scale of DeepSeek-LLM 1.3B, it may not possess sufficient capacity to thoroughly integrate code and mathematical data at the same time.

\subsubsection{ArXiv Papers Seem Ineffective in Improving Mathematical Reasoning}

ArXiv papers are commonly included as a component of math pre-training data \citep{minerva,gpt-f,llemma,mathpile}. However, detailed analysis regarding their impact on mathematical reasoning has not been extensively conducted. Perhaps counter-intuitively, according to our experiments, arXiv papers seem ineffective in improving mathematical reasoning. We experiment with models of different sizes, including DeepSeek-LLM 1.3B and DeepSeek-Coder-Base-v1.5 7B \citep{deepseek-coder}, using arXiv corpora that underwent varied processing pipelines:

\begin{itemize}[topsep=0pt]
    \item \textbf{MathPile} \citep{mathpile}: a dataset comprising 8.9 billion tokens, assembled through the application of heuristic cleaning and filtering procedures, with scientific papers from arXiv making up more than 85\% of the content;
    \item \textbf{ArXiv-RedPajama} \citep{redpajama}: a collection encompassing all arXiv LaTeX files, from which preambles, comments, macros, and bibliographies have been excised, amounting to 28.0 billion tokens in total.
\end{itemize}

In our study, we independently pre-train DeepSeek-LLM 1.3B on 150B tokens and DeepSeek-Coder-Base-v1.5 7B on 40B tokens for each arXiv dataset. Our findings indicate that utilizing arXiv papers alone does not lead to enhancements in mathematical reasoning abilities. When trained exclusively with arXiv data, neither model exhibits significant gains—and in some cases, performance even declines—across a range of mathematical evaluation sets of varying difficulty analyzed here. These assessment sets comprise quantitative reasoning tasks such as GSM8K and MATH (see Table \ref{tab:arxiv-cot}), multiple-choice evaluations like MMLU-STEM (Table \ref{tab:arxiv-cot}), as well as formal mathematics tests such as miniF2F (see Table \ref{tab:arxiv_atp}).

However, this conclusion has its limitations and should be taken with a grain of salt. We have not yet studied:

\begin{itemize}[topsep=0pt]
    \item How arXiv tokens influence additional mathematics-related applications outside the scope of this study, such as transforming formal theorem statements or proofs into informal language;
    \item The potential outcomes when arXiv tokens are integrated with alternative data sources;
    \item If increasing the model size could further reveal advantages associated with arXiv paper data.
\end{itemize}

Thus, further exploration is required, which we leave for future studies.

\subsection{Insights of Reinforcement Learning}
\subsubsection{Towards to a Unified Paradigm} This section introduces a unified framework for examining a variety of training strategies, including SFT, RFT, DPO, PPO, and GRPO. Additionally, we carry out empirical studies to investigate the key elements within this unified framework. In general, the gradient of a given training algorithm with respect to its parameter $\theta$ can be expressed as follows:

\begin{equation}
    \nabla_{\theta}\mathcal{J}_{\textcolor{red}{\mathcal{A}}}(\theta) = \mathbb{E}[\underbrace{(q,o) \sim \textcolor{red}{\mathcal{D}}}_{Data \ Source}]\left( \frac{1}{|o|} \sum_{t=1}^{|o|}  \underbrace{GC_{{\mathcal{A}}}(q, o, t, \textcolor{red}{\pi_{{rf}}})}_{Gradient \ Coefficient}  \nabla_{\theta}\log \pi_{\theta}(o_t | q, o_{<t})\right).
\label{eq:objective}
\end{equation}

There exist three key components: 1) \textit{Data Source $\mathcal{D}$}, which determines the training data; 2) \textit{Reward Function $\pi_{{rf}}$}, which is the source of the training reward signal; 3) \textit{Algorithm $\mathcal{A}$}: which processes the training data and the reward signal to the gradient coefficient $GC$ that determines the magnitude of the penalty or reinforcement for the data. We analyze several representative methods based on such a unified paradigm:

\begin{itemize}
    \item  \textbf{Supervised Fine-tuning (SFT)}: SFT involves adapting a pretrained model using curated SFT datasets selected by human annotators.
    \item \textbf{Rejection Sampling Fine-tuning (RFT)}: RFT builds upon the SFT model by performing an additional fine-tuning step. This step uses responses generated by the SFT model itself in response to SFT questions, but only retains those outputs that are verified as correct for further training.
    \item \textbf{Direct Preference Optimization (DPO)}: DPO enhances the SFT model by conducting fine-tuning on a dataset of augmented responses produced by the SFT model, optimizing with a pairwise DPO loss function.
    \item \textbf{Online Rejection Sampling Fine-tuning (Online RFT)}: In contrast to RFT, Online RFT starts with the SFT model as the policy model and continually updates it by fine-tuning on new, augmented responses that are sampled dynamically from the current version of the policy model.
    \item \textbf{PPO/GRPO}: PPO/GRPO begins by initializing the policy with the SFT model and subsequently applies reinforcement learning, utilizing outputs sampled on-the-fly from the evolving policy model for further updates.
\end{itemize}

Table \ref{tab:data_policy} provides an overview of the constituent elements of these approaches. For a comprehensive explanation of the derivation procedure, consult Appendix \ref{app:analysis-rl}.

\paragraph{Observation about Data Source} We categorize the sources of data into two main types: online sampling and offline sampling. Online sampling refers to data collected directly from the outcomes generated by the live, continuously learning policy model during training. In contrast, offline sampling involves gathering training data based on outputs produced by the pre-trained SFT model. Methods such as RFT and DPO employ the offline approach, whereas Online RFT and GRPO utilize the online approach.

As illustrated in Figure \ref{fig:rl-analysis}, our results indicate that Online RFT achieves markedly better performance than RFT across both evaluated benchmarks. In the initial phases of training, Online RFT and RFT yield similar outcomes, but as training progresses, Online RFT secures a clear performance lead, thereby highlighting the benefits of utilizing an online training framework. This outcome is expected: during the early training period, there is minimal divergence between the actor and the SFT model, resulting in sampled data that reflect only subtle distinctions. As training advances, however, the actor-generated samples increasingly diverge from those of the SFT model, and leveraging up-to-date samples through online data collection proves increasingly advantageous.

\paragraph{Observation about Gradient Coefficient} The mechanism by which the algorithm utilizes the reward signal to compute the gradient coefficient is pivotal for model parameter updates. In our study, we decompose the reward function into two components: `Rule' and `Model'. The `Rule' component evaluates response quality by verifying answer correctness, whereas the `Model' involves the training of a reward model tasked with scoring responses—this model is trained on datasets curated from rule-based judgments. A critical distinction between GRPO and Online RFT emerges in Equations \ref{eq:GC-RFT} and \ref{eq:GC-GRPO}: GRPO dynamically modifies its gradient coefficient as a function of the reward values generated by the reward model, thus enabling a nuanced approach that can either amplify or diminish the reinforcement signal based on the relative scores of different responses. Conversely, Online RFT does not possess this adaptive penalization mechanism; it offers uniform positive reinforcement to all correct responses and applies no penalization to incorrect ones, leading to equal intensity of updates for all correctly answered instances.

As illustrated in Figure \ref{fig:rl-analysis}, GRPO outperforms online RFT, underscoring the effectiveness of modifying gradient coefficients for positive and negative updates. Moreover, the results show that GRPO+PS achieves better outcomes than GRPO+OS, demonstrating the advantages conferred by incorporating step-aware, fine-grained gradient coefficients. Additionally, we investigate iterative RL by performing two sequential iterations in our experimental setup. Figure \ref{fig:iter-rl} reveals that applying iterative RL leads to marked performance gains, with the most pronounced improvement occurring during the initial iteration.

\subsubsection{Why RL Works?} In this study, we apply reinforcement learning using a selected portion of instruction tuning data and observe notable improvements over models trained solely with instruction tuning. To elucidate the effectiveness of reinforcement learning, we assess both the Pass@K and Maj@K accuracies for the Instruct and RL models across two distinct benchmarks. Figure \ref{fig:maj-pass} reveals that RL leads to better performance in terms of Maj@K accuracy, without a corresponding increase in Pass@K. These results imply that RL primarily strengthens the output distribution's resilience—\textbf{specifically, the gains appear to result from elevating the correct answer within the TopK outputs, rather than fundamentally augmenting the model’s underlying skills.} A parallel observation was made by \citep{wang2023making}, who identified a \textbf{misalignment problem} affecting reasoning performance in SFT models; they demonstrated that preference alignment strategies can enhance reasoning in such models \citep{yuan2023rrhf,song2023preference,wang2023making}.

\subsubsection{How to Achieve More Effective RL?}
\label{sec:effec_rl}
Our findings indicate that RL performs effectively in the context of mathematical reasoning challenges. Moreover, we introduce a cohesive framework for interpreting a range of prominent training approaches. This framework characterizes all strategies as instances of either direct or approximate RL methodologies. As outlined in Equation \ref{eq:objective}, the paradigm comprises three principal elements: Data Source, Algorithm, and Reward Function. We also discuss possible avenues for advancing each of these components in future work.

\paragraph{Data Source} The data source serves as the foundational input for all training strategies. Within RL, this term specifically denotes the unlabeled questions along with their associated outputs generated via sampling from the policy model. In this study, our approach restricts data sources to those questions originating from the instruction tuning phase, utilizing a basic nucleus sampling method for output generation. We speculate that this limitation may partly explain why our RL pipeline yields improvements solely in Maj@K metrics. Moving forward, we aim to extend our RL pipeline's application to question prompts outside the distribution of our current data, while also incorporating \textbf{advanced sampling (decoding) strategies}, such as tree-search-based methods \citep{yao2023tree}. In addition, employing \textbf{efficient inference techniques} \citep{xia-etal-2023-speculative,leviathan2023fast,kwon2023efficient,xia2024unlocking}, which influence the exploration effectiveness of policy models, is expected to have a substantial impact.

\paragraph{Algorithms} Algorithms utilize the available data and the associated reward signal to compute gradient coefficients, which subsequently adjust the model parameters. Drawing upon Equation \ref{eq:objective}, current approaches overwhelmingly place complete \textbf{TRUST} in the reward function's signal to modulate the conditional probability of specific tokens, either promoting or demoting them accordingly. Nonetheless, the assumption that the reward signal is consistently dependable does not always hold true, particularly in the context of highly intricate tasks. For instance, the PRM800K datasets \citep{lightman2023let}, despite being thoroughly labeled by experienced annotators, still exhibit an error rate of about 20\% in their annotations\footnote{\url{https://github.com/openai/prm800k/issues/12\#issuecomment-1728491852}}. Therefore, we aim to investigate reinforcement learning algorithms capable of maintaining robustness in the presence of noisy reward signals. Our expectation is that such \textbf{WEAK-TO-STRONG} \citep{burns2023weak} alignment methodologies could introduce significant paradigm shifts in how learning algorithms are designed.

\paragraph{Reward Function} The reward function serves as the principal source of learning signals during training. In reinforcement learning (RL), this function typically takes the form of a neural reward model. We identify three primary avenues for advancing reward models: 1) \textbf{Methods for improving the generalization capacity of reward models.} It is crucial that the reward model generalizes well to address out-of-distribution queries and more sophisticated decoding results; otherwise, RL may simply reinforce the existing distribution of LLMs instead of enhancing their core competencies; 2) \textbf{Approaches to represent reward model uncertainty.} Capturing uncertainty could play a pivotal role in bridging the limitations of underperforming reward models with algorithms that progress from weak to strong learners; 3) \textbf{Strategies for the efficient construction of high-quality process reward models,} which can deliver granular and informative feedback throughout the reasoning process \citep{lightman2023let,wang2023math}.

\section{Conclusion, Limitation, and Future Work}  
We introduce DeepSeekMath, a model that surpasses all publicly available systems on the MATH benchmark for competition-level mathematics and nears the results of proprietary models. Built by initializing from DeepSeek-Coder-v1.5 7B, DeepSeekMath undergoes extended training on 500B tokens, which notably includes 120B mathematics-related tokens extracted from Common Crawl. Our detailed ablation analyses reveal that data derived from web sources holds substantial promise for generating high-quality mathematical datasets, whereas data from arXiv appears less beneficial than anticipated. We propose Group Relative Policy Optimization (GRPO), an alternative to Proximal Policy Optimization (PPO), capable of enhancing mathematical reasoning performance while reducing memory demands. Experimental evidence indicates the effectiveness of GRPO, even for the high-performing DeepSeekMath-Instruct 7B model on benchmark tasks. Additionally, we present a unified framework to contextualize a range of reinforcement learning techniques and outline several promising future research directions for improved reinforcement learning in mathematical reasoning.

While DeepSeekMath~demonstrates strong performance on quantitative reasoning tasks, its abilities in geometry and theorem-proving lag behind those of proprietary models. For example, during our preliminary testing, the model was unable to solve questions involving triangles and ellipses, suggesting a possible bias in the types of data used during pre-training and fine-tuning stages. Moreover, due to the limitation in model size, DeepSeekMath~exhibits inferior few-shot learning capabilities when compared to GPT-4. Unlike GPT-4, which can enhance its results with few-shot prompts, DeepSeekMath~shows nearly identical outcomes in both zero-shot and few-shot settings. Moving forward, we intend to refine our engineered data selection process to assemble a more robust and diverse pre-training dataset. Additionally, we plan to investigate new strategies for enhancing LLM reinforcement learning effectiveness, as discussed in Section \ref{sec:effec_rl}.

\end{document}